% GENERATED FILE. Edit canonical lessons, then run npm run generate:books.
% canonical-source-hash: fnv1a64:54d1884321f25991
% canonical-lessons: TA-C513-tolilmunaivor, TA-C513-alocakar, TA-C513-kollar, TA-C513-necavalar, TA-C513-panippen, TA-R513-review-506-513

\chapter{People at Work}
\label{ch:ta-people-at-work}
\hlchaptermodality{513}

\begin{chapteropening}
\textbf{By the end of this chapter:} \emph{I can name more people by the work they do.}

Complete the last lesson of chapter 513: I can name more people by the work they do.
\end{chapteropening}

\section[\texorpdfstring{\ta{தொழில்முனைவோர்} (toḻilmuṉaivōr)}{தொழில்முனைவோர் (toḻilmuṉaivōr)}]{\ta{தொழில்முனைவோர்} (toḻilmuṉaivōr) — an entrepreneur}

\label{lesson:TA-C513-tolilmunaivor}



\begin{quote}
Before the new one: say the Tamil for a transfer, then the Tamil for to start.
\end{quote}

\subsection*{You'll want to know: \ta{தொழில்முனைவோர்}}
\textbf{\ta{தொழில்முனைவோர்}} — \emph{toḻilmuṉaivōr} — \textquotedblleft{}an entrepreneur\textquotedblright{}.

\begin{quote}
\ta{இளம்} \ta{தொழில்முனைவோர்} \ta{பலர்} \ta{இந்த} \ta{நகரத்தில்} \ta{இருக்கிறார்கள்}.
\end{quote}

\begin{quote}
\emph{iḷam toḻilmuṉaivōr palar inta nakarattil irukkiṟārkaḷ.}
\end{quote}

\begin{quote}
There are many young entrepreneurs in this city.
\end{quote}

The ending -\ta{ஓர்} makes it a word for people; it serves for one person and for many.

\begin{cousinweb}
\textbf{\ta{தொழில்முனைவோர்}} (\emph{toḻilmuṉaivōr}) is a modern native compound: \ta{தொழில்} (\emph{toḻil}), \textquotedblleft{}a trade, work\textquotedblright{} (DEDR 3524), and \ta{முனைவோர்}, \textquotedblleft{}those who set out to do something\textquotedblright{}, from \ta{முனை} (\emph{muṉai}), \textquotedblleft{}to strive\textquotedblright{}.
\end{cousinweb}

\begin{grammarlens}[title={What you've built}]
One more word for describing what you lived and felt.
\end{grammarlens}

\subsection*{Your turn}
\begin{itemize}
\raggedright
  \item \emph{Say it:} \emph{toḻilmuṉaivōr}
  \item \emph{Say it:} \emph{toḻilmuṉaivōr}, once more
  \item \emph{Say it:} \emph{iḷam toḻilmuṉaivōr palar inta nakarattil irukkiṟārkaḷ.} — There are many young entrepreneurs in this city.
  \item \emph{Recall:} say the Tamil for a transfer, then the Tamil for to start, then say \emph{toḻilmuṉaivōr} again
\end{itemize}

\begin{tcolorbox}[breakable,colback=teal!4,colframe=teal!35!black,title={Before you move on}]
What does \ta{தொழில்முனைவோர்} mean? (\textquotedblleft{}An entrepreneur\textquotedblright{}.) Say it once more.
\end{tcolorbox}
\section[\texorpdfstring{\ta{ஆலோசகர்} (ālōcakar)}{ஆலோசகர் (ālōcakar)}]{\ta{ஆலோசகர்} (ālōcakar) — an adviser; a consultant}

\label{lesson:TA-C513-alocakar}



\begin{quote}
Before the new one: say the Tamil for to start, then the Tamil for an entrepreneur.
\end{quote}

\subsection*{You'll want to know: \ta{ஆலோசகர்}}
\textbf{\ta{ஆலோசகர்}} — \emph{ālōcakar} — \textquotedblleft{}an adviser; a consultant\textquotedblright{}.

\begin{quote}
\ta{நிறுவனம்} \ta{ஒரு} \ta{நிதி} \ta{ஆலோசகரை} \ta{அமர்த்தியது}.
\end{quote}

\begin{quote}
\emph{niṟuvaṉam oru niti ālōcakarai amarttiyatu.}
\end{quote}

\begin{quote}
The company hired a financial adviser.
\end{quote}

Built on \ta{ஆலோசனை}, advice. The person advised takes -\ta{க்கு}: \ta{அரசுக்கு} \ta{ஆலோசகர்}, an adviser to the government.

\begin{cousinweb}
\textbf{\ta{ஆலோசகர்}} (\emph{ālōcakar}) is borrowed from Sanskrit: the stem of \emph{ālocana}, \textquotedblleft{}consideration, reflection\textquotedblright{} (Monier-Williams), with the Tamil personal ending -\ta{அர்}.
\end{cousinweb}

\begin{grammarlens}[title={What you've built}]
One more word for describing what you lived and felt.
\end{grammarlens}

\subsection*{Your turn}
\begin{itemize}
\raggedright
  \item \emph{Say it:} \emph{ālōcakar}
  \item \emph{Say it:} \emph{ālōcakar}, once more
  \item \emph{Say it:} \emph{niṟuvaṉam oru niti ālōcakarai amarttiyatu.} — The company hired a financial adviser.
  \item \emph{Recall:} say the Tamil for to start, then the Tamil for an entrepreneur, then say \emph{ālōcakar} again
\end{itemize}

\begin{tcolorbox}[breakable,colback=teal!4,colframe=teal!35!black,title={Before you move on}]
What does \ta{ஆலோசகர்} mean? (\textquotedblleft{}An adviser; a consultant\textquotedblright{}.) Say it once more.
\end{tcolorbox}
\section[\texorpdfstring{\ta{கொல்லர்} (kollar)}{கொல்லர் (kollar)}]{\ta{கொல்லர்} (kollar) — a blacksmith}

\label{lesson:TA-C513-kollar}



\begin{quote}
Before the new one: say the Tamil for an entrepreneur, then the Tamil for an adviser.
\end{quote}

\subsection*{You'll want to know: \ta{கொல்லர்}}
\textbf{\ta{கொல்லர்}} — \emph{kollar} — \textquotedblleft{}a blacksmith\textquotedblright{}.

\begin{quote}
\ta{கொல்லர்} \ta{பழுத்த} \ta{இரும்பை} \ta{அடித்து} \ta{ஒரு} \ta{கத்தி} \ta{செய்தார்}.
\end{quote}

\begin{quote}
\emph{kollar paḻutta irumpai aṭittu oru katti ceytār.}
\end{quote}

\begin{quote}
The blacksmith hammered the red-hot iron into a knife.
\end{quote}

A word for one person given with respect; the older form is \ta{கொல்லன்}.

\begin{cousinweb}
\textbf{\ta{கொல்லர்}} (\emph{kollar}) is native Dravidian: DEDR 2133 gives \ta{கொல்} (\emph{kol}), \textquotedblleft{}working in iron\textquotedblright{}, and \ta{கொல்லன்} (\emph{kollaṉ}), \textquotedblleft{}a blacksmith\textquotedblright{}; \ta{கொல்லர்} is the form with the respectful -\ta{அர்}.
\end{cousinweb}

\begin{grammarlens}[title={What you've built}]
One more word for describing what you lived and felt.
\end{grammarlens}

\subsection*{Your turn}
\begin{itemize}
\raggedright
  \item \emph{Say it:} \emph{kollar}
  \item \emph{Say it:} \emph{kollar}, once more
  \item \emph{Say it:} \emph{kollar paḻutta irumpai aṭittu oru katti ceytār.} — The blacksmith hammered the red-hot iron into a knife.
  \item \emph{Recall:} say the Tamil for an entrepreneur, then the Tamil for an adviser, then say \emph{kollar} again
\end{itemize}

\begin{tcolorbox}[breakable,colback=teal!4,colframe=teal!35!black,title={Before you move on}]
What does \ta{கொல்லர்} mean? (\textquotedblleft{}A blacksmith\textquotedblright{}.) Say it once more.
\end{tcolorbox}
\section[\texorpdfstring{\ta{நெசவாளர்} (necavāḷar)}{நெசவாளர் (necavāḷar)}]{\ta{நெசவாளர்} (necavāḷar) — a weaver}

\label{lesson:TA-C513-necavalar}



\begin{quote}
Before the new one: say the Tamil for an adviser, then the Tamil for a blacksmith.
\end{quote}

\subsection*{You'll want to know: \ta{நெசவாளர்}}
\textbf{\ta{நெசவாளர்}} — \emph{necavāḷar} — \textquotedblleft{}a weaver\textquotedblright{}.

\begin{quote}
\ta{எங்கள்} \ta{கிராமத்தில்} \ta{பல} \ta{நெசவாளர்கள்} \ta{சேலைகள்} \ta{நெய்கிறார்கள்}.
\end{quote}

\begin{quote}
\emph{eṅkaḷ kirāmattil pala necavāḷarkaḷ cēlaikaḷ neykiṟārkaḷ.}
\end{quote}

\begin{quote}
In our village many weavers weave saris.
\end{quote}

Plural: \ta{நெசவாளர்கள்}, \emph{necavāḷarkaḷ}.

The ending -\ta{ஆளர்} makes \textquotedblleft{}one who …\textquotedblright{}; the verb to weave is \ta{நெய்} (\emph{ney}), which also names ghee.

\begin{cousinweb}
\textbf{\ta{நெசவாளர்}} (\emph{necavāḷar}) is native Dravidian: DEDR 3745 gives \ta{நெய்} (\emph{ney}), \textquotedblleft{}to weave\textquotedblright{}, and \ta{நெசவு} (\emph{necavu}), \textquotedblleft{}weaving\textquotedblright{}; -\ta{ஆளர்} makes \textquotedblleft{}one who weaves\textquotedblright{}.
\end{cousinweb}

\begin{grammarlens}[title={What you've built}]
One more word for describing what you lived and felt.
\end{grammarlens}

\subsection*{Your turn}
\begin{itemize}
\raggedright
  \item \emph{Say it:} \emph{necavāḷar}
  \item \emph{Say it:} \emph{necavāḷar}, once more
  \item \emph{Say it:} \emph{eṅkaḷ kirāmattil pala necavāḷarkaḷ cēlaikaḷ neykiṟārkaḷ.} — In our village many weavers weave saris.
  \item \emph{Recall:} say the Tamil for an adviser, then the Tamil for a blacksmith, then say \emph{necavāḷar} again
\end{itemize}

\begin{tcolorbox}[breakable,colback=teal!4,colframe=teal!35!black,title={Before you move on}]
What does \ta{நெசவாளர்} mean? (\textquotedblleft{}A weaver\textquotedblright{}.) Say it once more.
\end{tcolorbox}
\section[\texorpdfstring{\ta{பணிப்பெண்} (paṇippeṇ)}{பணிப்பெண் (paṇippeṇ)}]{\ta{பணிப்பெண்} (paṇippeṇ) — a maid; a housekeeper}

\label{lesson:TA-C513-panippen}



\begin{quote}
Before the new one: say the Tamil for a blacksmith, then the Tamil for a weaver.
\end{quote}

\subsection*{You'll want to know: \ta{பணிப்பெண்}}
\textbf{\ta{பணிப்பெண்}} — \emph{paṇippeṇ} — \textquotedblleft{}a maid; a housekeeper\textquotedblright{}.

\begin{quote}
\ta{அந்த} \ta{வீட்டில்} \ta{ஒரு} \ta{பணிப்பெண்} \ta{சமையல்} \ta{செய்கிறாள்}.
\end{quote}

\begin{quote}
\emph{anta vīṭṭil oru paṇippeṇ camaiyal ceykiṟāḷ.}
\end{quote}

\begin{quote}
In that house a maid does the cooking.
\end{quote}

\ta{பணி} \textquotedblleft{}service\textquotedblright{} + \ta{பெண்} \textquotedblleft{}a woman\textquotedblright{}. A man in the same work is a \ta{பணியாளர்}.

\begin{cousinweb}
\textbf{\ta{பணிப்பெண்}} (\emph{paṇippeṇ}) is a native compound: \ta{பணி} (\emph{paṇi}), \textquotedblleft{}service\textquotedblright{} (DEDR 3884), and \ta{பெண்} (\emph{peṇ}), \textquotedblleft{}a woman\textquotedblright{} (DEDR 4395).
\end{cousinweb}

\begin{grammarlens}[title={What you've built}]
One more word for describing what you lived and felt.
\end{grammarlens}

\subsection*{Your turn}
\begin{itemize}
\raggedright
  \item \emph{Say it:} \emph{paṇippeṇ}
  \item \emph{Say it:} \emph{paṇippeṇ}, once more
  \item \emph{Say it:} \emph{anta vīṭṭil oru paṇippeṇ camaiyal ceykiṟāḷ.} — In that house a maid does the cooking.
  \item \emph{Recall:} say the Tamil for a blacksmith, then the Tamil for a weaver, then say \emph{paṇippeṇ} again
\end{itemize}

\begin{tcolorbox}[breakable,colback=teal!4,colframe=teal!35!black,title={Before you move on}]
What does \ta{பணிப்பெண்} mean? (\textquotedblleft{}A maid; a housekeeper\textquotedblright{}.) Say it once more.
\end{tcolorbox}
\section[(dialogue)]{Looking back at work and careers}

\label{lesson:TA-R513-review-506-513}



\begin{quote}
Say the Tamil for a salaried job and a CV.
\end{quote}

\subsection*{Your turn}
Say these aloud:

\begin{itemize}
\raggedright
  \item a salaried job, a CV, to hire, full-time, an assistant
  \item a workplace, to be employed, a secretary, a file, to supervise
  \item wages, leave from work, overtime, an accountant, a clerk
  \item a promotion, an ambition, to work hard, to encourage, an expert
  \item production, to fit, to coordinate, to stack, to take measurements
  \item exports, imports, business, a businessman, a partner
  \item dismissal, to resign, to retire, a transfer, to start
  \item an entrepreneur, an adviser, a blacksmith, a weaver, a maid
\end{itemize}

\begin{tcolorbox}[breakable,colback=teal!4,colframe=teal!35!black,title={Before you move on}]
A salaried job? (\textbf{\ta{உத்தியோகம்}}, \emph{uttiyōkam}.) A CV? (\textbf{\ta{சுயவிவரம்}}, \emph{cuyavivaram}.) To hire? (\textbf{\ta{அமர்த்து}}, \emph{amarttu}.) Full-time? (\textbf{\ta{முழுநேர}}, \emph{muḻunēra}.) An assistant? (\textbf{\ta{உதவியாளர்}}, \emph{utaviyāḷar}.) A workplace? (\textbf{\ta{பணியிடம்}}, \emph{paṇiyiṭam}.) To be employed? (\textbf{\ta{பணிபுரி}}, \emph{paṇipuri}.) A secretary? (\textbf{\ta{செயலாளர்}}, \emph{ceyalāḷar}.) A file? (\textbf{\ta{கோப்பு}}, \emph{kōppu}.) To supervise? (\textbf{\ta{மேற்பார்வையிடு}}, \emph{mēṟpārvaiyiṭu}.) Wages? (\textbf{\ta{ஊதியம்}}, \emph{ūtiyam}.) Leave from work? (\textbf{\ta{விடுப்பு}}, \emph{viṭuppu}.) Overtime? (\textbf{\ta{கூடுதல்} \ta{நேரம்}}, \emph{kūṭutal nēram}.) An accountant? (\textbf{\ta{கணக்காளர்}}, \emph{kaṇakkāḷar}.) A clerk? (\textbf{\ta{எழுத்தர்}}, \emph{eḻuttar}.) A promotion? (\textbf{\ta{பதவி} \ta{உயர்வு}}, \emph{patavi uyarvu}.) An ambition? (\textbf{\ta{லட்சியம்}}, \emph{laṭciyam}.) To work hard? (\textbf{\ta{உழை}}, \emph{uḻai}.) To encourage? (\textbf{\ta{ஊக்குவி}}, \emph{ūkkuvi}.) An expert? (\textbf{\ta{நிபுணர்}}, \emph{nipuṇar}.) Production? (\textbf{\ta{உற்பத்தி}}, \emph{uṟpatti}.) To fit? (\textbf{\ta{பொருத்து}}, \emph{poruttu}.) To coordinate? (\textbf{\ta{ஒருங்கிணை}}, \emph{oruṅkiṇai}.) To stack? (\textbf{\ta{அடுக்கு}}, \emph{aṭukku}.) To take measurements? (\textbf{\ta{அள}}, \emph{aḷa}.) Exports? (\textbf{\ta{ஏற்றுமதி}}, \emph{ēṟṟumati}.) Imports? (\textbf{\ta{இறக்குமதி}}, \emph{iṟakkumati}.) Business? (\textbf{\ta{வணிகம்}}, \emph{vaṇikam}.) A businessman? (\textbf{\ta{தொழிலதிபர்}}, \emph{toḻilatipar}.) A partner? (\textbf{\ta{கூட்டாளி}}, \emph{kūṭṭāḷi}.) Dismissal? (\textbf{\ta{பணிநீக்கம்}}, \emph{paṇinīkkam}.) To resign? (\textbf{\ta{பதவி} \ta{விலகு}}, \emph{patavi vilaku}.) To retire? (\textbf{\ta{ஓய்வு} \ta{பெறு}}, \emph{ōyvu peṟu}.) A transfer? (\textbf{\ta{இடமாற்றம்}}, \emph{iṭamāṟṟam}.) To start? (\textbf{\ta{தொடங்கு}}, \emph{toṭaṅku}.) An entrepreneur? (\textbf{\ta{தொழில்முனைவோர்}}, \emph{toḻilmuṉaivōr}.) An adviser? (\textbf{\ta{ஆலோசகர்}}, \emph{ālōcakar}.) A blacksmith? (\textbf{\ta{கொல்லர்}}, \emph{kollar}.) A weaver? (\textbf{\ta{நெசவாளர்}}, \emph{necavāḷar}.) A maid? (\textbf{\ta{பணிப்பெண்}}, \emph{paṇippeṇ}.)
\end{tcolorbox}
